% Armijo backtracking on the common quadratic at x0=(2,2)
% Generated by tikz/generate_c3.py; edit the generator and rerun.
\begin{tikzpicture}[font=\fontsize{8}{10}\selectfont,
  note/.style={font=\FigureNoteFont\fontsize{7.5}{9}\selectfont}]
\pgfplotsset{every axis/.append style={
  scale only axis, axis x line*=bottom, axis y line*=left,
  axis line style={FigureStroke,line width=.75pt}, tick style={FigureStroke,line width=.5pt},
  tick label style={font=\fontsize{7.5}{9}\selectfont,text=FigureInk},
  label style={font=\fontsize{8}{10}\selectfont,text=FigureInk},
  title style={font=\fontsize{8}{10}\selectfont,text=FigureInk,yshift=-1mm},
  legend style={font=\fontsize{7.5}{9}\selectfont,draw=none,fill=white,fill opacity=.9,text opacity=1,
    inner sep=1pt,row sep=-1pt},
  legend cell align=left}}
\begin{axis}[width=86mm,height=44mm,xmin=0,xmax=.15,ymin=0,ymax=95,
  xlabel={步长 $\eta$},ylabel={目标值},legend pos=north east,
  xtick={0,.03125,.0625,.09375,.125},xticklabels={$0$,$\frac1{32}$,$\frac1{16}$,$\frac3{32}$,$\frac18$}]
\fill[FigureGreenFill] (axis cs:0,0) rectangle (axis cs:0.09409,95);
\addplot[FigureInk,line width=1.1pt] coordinates {
    (0,34) (0.001,32.98) (0.002,31.977) (0.003,30.99) (0.004,30.019) (0.005,29.065)
    (0.006,28.127) (0.007,27.206) (0.008,26.3) (0.009,25.412) (0.01,24.539) (0.011,23.683)
    (0.012,22.844) (0.013,22.021) (0.014,21.214) (0.015,20.424) (0.016,19.65) (0.017,18.892)
    (0.018,18.151) (0.019,17.426) (0.02,16.718) (0.021,16.026) (0.022,15.35) (0.023,14.691)
    (0.024,14.048) (0.025,13.421) (0.026,12.811) (0.027,12.217) (0.028,11.64) (0.029,11.079)
    (0.03,10.535) (0.031,10.006) (0.032,9.4947) (0.033,8.9993) (0.034,8.5203) (0.035,8.0576)
    (0.036,7.6114) (0.037,7.1816) (0.038,6.7681) (0.039,6.3711) (0.04,5.9904) (0.041,5.6261)
    (0.042,5.2782) (0.043,4.9467) (0.044,4.6316) (0.045,4.3329) (0.046,4.0505) (0.047,3.7845)
    (0.048,3.535) (0.049,3.3018) (0.05,3.085) (0.051,2.8846) (0.052,2.7006) (0.053,2.5329)
    (0.054,2.3817) (0.055,2.2468) (0.056,2.1284) (0.057,2.0263) (0.058,1.9406) (0.059,1.8713)
    (0.06,1.8184) (0.061,1.7819) (0.062,1.7617) (0.063,1.758) (0.064,1.7706) (0.065,1.7997)
    (0.066,1.8451) (0.067,1.9069) (0.068,1.9851) (0.069,2.0796) (0.07,2.1906) (0.071,2.318)
    (0.072,2.4617) (0.073,2.6218) (0.074,2.7983) (0.075,2.9912) (0.076,3.2005) (0.077,3.4262)
    (0.078,3.6683) (0.079,3.9268) (0.08,4.2016) (0.081,4.4928) (0.082,4.8005) (0.083,5.1245)
    (0.084,5.4649) (0.085,5.8217) (0.086,6.1948) (0.087,6.5844) (0.088,6.9903) (0.089,7.4127)
    (0.09,7.8514) (0.091,8.3065) (0.092,8.778) (0.093,9.2659) (0.094,9.7702) (0.095,10.291)
    (0.096,10.828) (0.097,11.381) (0.098,11.951) (0.099,12.537) (0.1,13.14) (0.101,13.759)
    (0.102,14.394) (0.103,15.046) (0.104,15.714) (0.105,16.399) (0.106,17.1) (0.107,17.817)
    (0.108,18.551) (0.109,19.301) (0.11,20.067) (0.111,20.85) (0.112,21.65) (0.113,22.465)
    (0.114,23.297) (0.115,24.146) (0.116,25.01) (0.117,25.892) (0.118,26.789) (0.119,27.703)
    (0.12,28.634) (0.121,29.58) (0.122,30.543) (0.123,31.523) (0.124,32.519) (0.125,33.531)
    (0.126,34.56) (0.127,35.605) (0.128,36.666) (0.129,37.744) (0.13,38.839) (0.131,39.949)
    (0.132,41.076) (0.133,42.22) (0.134,43.379) (0.135,44.556) (0.136,45.748) (0.137,46.957)
    (0.138,48.183) (0.139,49.424) (0.14,50.682) (0.141,51.957) (0.142,53.248) (0.143,54.555)
    (0.144,55.879) (0.145,57.219) (0.146,58.575) (0.147,59.948) (0.148,61.337) (0.149,62.743)
    (0.15,64.165)};
\addlegendentry{$\varphi(\eta)=f(x-\eta\nabla f(x))$}
\addplot[FigureMuted,dotted,line width=.8pt] coordinates {
    (0,34) (0.033074,0)};
\addlegendentry{切线 $34-1028\eta$}
\addplot[FigureBlue,dashed,line width=.9pt] coordinates {
    (0,34) (0.15,-4.55)};
\addlegendentry{Armijo 线 $34-257\eta$}
\addplot[forget plot,only marks,mark=*,mark size=1.8pt,FigureRed] coordinates {(.125,33.531)};
\addplot[forget plot,only marks,mark=*,mark size=1.8pt,FigureGreen] coordinates {(.0625,1.7578)};
\node[note,text=FigureRed,anchor=south] at (axis cs:.125,35.531) {拒绝};
\node[note,text=FigureGreen,anchor=south] at (axis cs:.0625,3.7578) {接受};
\node[note,text=FigureGreen!70!FigureInk,anchor=north west] at (axis cs:.002,93) {充分下降区间};
\end{axis}
\end{tikzpicture}
